\section{Code Audit --- One-Line Comment
Invariant}\label{code-audit-one-line-comment-invariant}

\subsection{Audit Procedure}\label{audit-procedure}

Every Python file in \texttt{src/} is checked by
\texttt{scripts/check\_comments.py}, which: 1. Parses the file with
\texttt{ast.parse} to enumerate all \texttt{FunctionDef} and
\texttt{AsyncFunctionDef} nodes. 2. Tokenizes the file to locate
\texttt{COMMENT} tokens. 3. For each \texttt{def}, walks backward from
its line number through any intervening decorator lines (\texttt{@}) and
blank lines. A one-line comment must appear as a direct predecessor
(only decorators and blank lines allowed between the comment and the
\texttt{def}).

The invariant is from \texttt{CLAUDE.md} §5: ``Every function has a
one-line comment directly above its \texttt{def} line.''

\subsection{Audit Results}\label{audit-results}

\subsubsection{Phase 3 completion
(2026-04-18)}\label{phase-3-completion-2026-04-18}

\begin{verbatim}
python3 scripts/check_comments.py --path src/
OK — all defs have a one-line comment directly above. (src)
\end{verbatim}

\textbf{Files checked (26):} - \texttt{src/attacks/\_\_init\_\_.py} -
\texttt{src/attacks/inversion.py} - \texttt{src/attacks/membership.py} -
\texttt{src/attacks/similarity.py} - \texttt{src/attacks/utility.py} -
\texttt{src/attacks/verbatim.py} - \texttt{src/data/\_\_init\_\_.py} -
\texttt{src/data/\_fakers.py} - \texttt{src/data/\_vocab.py} -
\texttt{src/data/annotator.py} - \texttt{src/data/canary.py} -
\texttt{src/data/schemas.py} -
\texttt{src/data/synthetic\_monitoring.py} -
\texttt{src/data/synthetic\_protocol.py} -
\texttt{src/data/synthetic\_sae.py} -
\texttt{src/data/synthetic\_writing.py} -
\texttt{src/ngsp/\_\_init\_\_.py} - \texttt{src/ngsp/answer\_applier.py}
- \texttt{src/ngsp/dp\_mechanism.py} -
\texttt{src/ngsp/entity\_extractor.py} -
\texttt{src/ngsp/local\_model.py} - \texttt{src/ngsp/pipeline.py} -
\texttt{src/ngsp/proxy\_decoder.py} -
\texttt{src/ngsp/query\_synthesizer.py} -
\texttt{src/ngsp/remote\_client.py} - \texttt{src/ngsp/router.py} -
\texttt{src/ngsp/safe\_harbor.py}

\textbf{Violations:} 0

\subsection{Special Cases Handled}\label{special-cases-handled}

\textbf{Pydantic validators with decorator chains.} CLAUDE.md requires
the comment directly above the \texttt{def}, which means it must appear
between the last decorator and the \texttt{def} line (not above the
decorator chain). For example:

\begin{Shaded}
\begin{Highlighting}[]
\AttributeTok{@field\_validator}\NormalTok{(}\StringTok{"model\_id"}\NormalTok{)}
\AttributeTok{@classmethod}
\CommentTok{\# Reject empty model IDs early; from\_pretrained would give a confusing error later.}
\KeywordTok{def}\NormalTok{ \_nonempty\_model\_id(cls, v: }\BuiltInTok{str}\NormalTok{) }\OperatorTok{{-}\textgreater{}} \BuiltInTok{str}\NormalTok{:}
\NormalTok{    ...}
\end{Highlighting}
\end{Shaded}

The checker correctly handles this case: it walks backward from the
\texttt{def} line through decorator lines and blank lines, accepting a
comment at any point in that span.

\textbf{\texttt{\_\_init\_\_} files.} Files with no function definitions
(e.g., \texttt{src/ngsp/\_\_init\_\_.py},
\texttt{src/attacks/\_\_init\_\_.py}) pass trivially with zero defs and
zero violations.

\subsection{CI Integration}\label{ci-integration}

To run the check in CI, add to your test runner:

\begin{Shaded}
\begin{Highlighting}[]
\ExtensionTok{python3}\NormalTok{ scripts/check\_comments.py }\AttributeTok{{-}{-}path}\NormalTok{ src/}
\end{Highlighting}
\end{Shaded}

This exits with code 0 on success, code 1 on any violation. It can be
added as a pre-commit hook or a GitHub Actions step.
